\subsection{形式幂级数的泰勒公式}

设 \(\mathbf{X} = (\mathbf{X}_i)_{i\in I}\) 与 \(\mathbf{Y} = (\mathbf{Y}_i)_{i\in I}\) 是关于同一指标集 I 的两个未定元族。我们用 \(\mathbf{X} + \mathbf{Y}\) 表示族 \((\mathbf{X}_i + \mathbf{Y}_i)_{i\in I}\) ，它由 \(\mathbf{A}[[\mathbf{X},\mathbf{Y}]]\) 中的形式幂级数组成。显然，我们可以以 \(\mathbf{X}_i + \mathbf{Y}_i\) 代替 \(\mathbf{X}_i\) ，代入形式幂级数 \(u\in \mathbf{A}[[\mathbf{X}]]\) 中，结果记作 \(u(\mathbf{X} + \mathbf{Y})\) 。对每个 \(\nu \in \mathbf{N}^{(I)}\) ，我们用 \(\Delta^{\nu}u\) 表示 \(\mathbf{Y}^{\nu}\) 在形式幂级数 \(u(\mathbf{X} + \mathbf{Y})\) （视为属于 \(\mathbf{A}[[\mathbf{X}]][[\mathbf{Y}]]\) ）中的系数（III，第456页）。换言之，我们有

\[
u (\mathbf {X} + \mathbf {Y}) = \sum_ {\nu} \Delta^ {\nu} u (\mathbf {X}) \cdot \mathbf {Y} ^ {\nu} \quad (u \in \mathrm{A} [ [ \mathbf {X} ] ]).\tag{6}
\]

以 (0, X) 代替 (X, Y) 作代入，我们得到

\[
u (\mathbf {X}) = \sum_ {\nu} \Delta^ {\nu} u (0) \cdot \mathbf {X} ^ {\nu};\tag{7}
\]

换言之，\(\Delta^{\nu}u\) 的常数项是 \(\mathbf{X}^{\nu}\) 在 \(u\) 中的系数。由于映射 \(u \mapsto u(\mathbf{X} + \mathbf{Y})\) （从 \(\mathbf{A}[[\mathbf{X}]]\) 到 \(\mathbf{A}[[\mathbf{X}, \mathbf{Y}]]\) ）是连续的，映射 \(u \mapsto \Delta^{\nu}u\) （从 \(\mathbf{A}[[\mathbf{X}]]\) 到其自身）也是连续的。

如同多项式的情形（IV，第7页），我们可以证明公式

\begin{enumerate}
\def\labelenumi{(\arabic{enumi})}
\setcounter{enumi}{7}
\tightlist
\item
\end{enumerate}

\[
\Delta^ {\sigma} (u v) = \sum_ {\nu + \rho = \sigma} \Delta^ {\nu} (u) \Delta^ {\rho} (v),
\]

\[
\Delta^ {\rho} \Delta^ {\sigma} u = \frac {(\rho + \sigma) !}{\rho ! \sigma !} \Delta^ {\rho + \sigma} u.\tag{9}
\]

二项式公式（I，第99页，推论 2）给出 \(\Delta^{\nu}u\) 当 \(u = \sum_{\lambda} \alpha_{\lambda} X^{\lambda}\) 时的如下值

\[
\Delta^ {\nu} u = \sum_ {\lambda} \alpha_ {\lambda + \nu} \frac {(\lambda + \nu) !}{\lambda ! \nu !} X ^ {\lambda}.\tag{10}
\]

特别地考虑情形 \(\nu = \varepsilon_{i}\) ，即 \(\nu_{i} = 1\) ，\(\nu_{j} = 0\) 对 \(j \neq i\) 。我们将记 \(D_{i}u = \Delta^{\varepsilon_{i}}u\) ；换言之，\(D_{i}u\) 是 \(Y_{i}\) 在 \(u(\mathbf{X} + \mathbf{Y})\) 中的系数。于是由 (10) 我们有

\[
\mathrm{D} _ {i} u = \sum_ {\lambda} (\lambda_ {i} + 1) \alpha_ {\lambda + \varepsilon_ {i}} X ^ {\lambda};\tag{11}
\]

特别地，我们有 \(\mathbf{D}_{i}(\mathbf{X}_{i})=1\) 与 \(\mathbf{D}_{i}(\mathbf{X}_{j})=0\) 对 \(j\neq i\) 。公式 (8) 表明 \(D_{i}\) 是 A{[}{[}X{]}{]} 的一个导子，并且由 (9) 我们推出关系

\[
\mathbf {D} ^ {\nu} \boldsymbol {u} = \nu ! \Delta^ {\nu} \boldsymbol {u}\tag{12}
\]

如同多项式的情形（IV，第8页）（我们已置 \(\mathbf{D}^{\nu} = \prod_{i\in I}\mathbf{D}_{i}^{\nu_{i}}\) ，对 \(\nu = (\nu_{i})_{i\in I}\) 属于 \(\mathbf{N}^{(I)}\) ）。当 \(\mathbf{A}\) 是一个 \(\mathbf{Q}\) -代数时，公式 (6)、(7) 与 (12) 蕴含「泰勒公式」：

\begin{enumerate}
\def\labelenumi{(\arabic{enumi})}
\setcounter{enumi}{12}
\tightlist
\item
\end{enumerate}

\[
u (\mathbf {X} + \mathbf {Y}) = \sum_ {\nu} \frac {1}{\nu !} D ^ {\nu} u (\mathbf {X}). \mathbf {Y} ^ {\nu},
\]

\[
u (\mathbf {X}) = \sum_ {\nu} \frac {1}{\nu !} D ^ {\nu} u (0) \cdot X ^ {\nu}.\tag{14}
\]

注记. --- 1) 我们常说 \(D_i u\) 是 \(u\) 关于 \(X_i\) 的偏导数；我们也使用记号 \(D_{X_i} u, \frac{\partial u}{\partial X_i}\) 与 \(u_{X_i}'\) 。对单个未定元 \(X\) ，唯一的偏导数 \(Du\) （也写作 \(\frac{du}{dX}\) 或 \(u'\) ）称为 \(u\) 的导数。

\begin{enumerate}
\def\labelenumi{\arabic{enumi})}
\setcounter{enumi}{1}
\item
  公式 (9) 表明自同态 \(\Delta^{\rho}\) （即 A-模 \(\mathbf{A}[[\mathbf{X}]]\) 的自同态）两两交换。因此对自同态 \(D_i\) 而言同样成立。
\item
  如果 \(u \in \mathbf{A}[(\mathbf{X}_i)_{i \in I}]\) 是多项式，则在 IV 第6页与第7页定义的多项式 \(\Delta^{\rho} u\) 与 \(D_i u\) 跟用同样符号表示的形式幂级数一致。
\end{enumerate}

